\clearpage
\chapter*{Đính chính và bổ sung cho Báo cáo số~236}
\addcontentsline{toc}{chapter}{Đính chính và bổ sung cho Báo cáo số~236}
\sid{e236}
\markboth{Đính chính 236}{Đính chính 236}
\src{303}

\noindent\textbf{8. TDTE VI: Các lược đồ Picard. Các tính chất tổng quát.}

\begin{center}
[Séminaire Bourbaki, tập 14, 1961/62, số 236, 23 trang.]
\end{center}

\noindent\textbf{Trang 12, dòng 16.} --- Thay ``đường thẳng afin'' bằng
``đường thẳng afin bỏ gốc tọa độ''.

\noindent\textbf{Dòng 18.} --- Thay ``$X[t]$'' bằng
``$X[t,t^{-1}]$''.

\noindent\textbf{Trang 16, Mệnh đề 3.1, dòng thứ 6.} --- Thay ``Bây giờ giả
thiết ...'' bằng ``Bây giờ, từ giả thiết ...''.

\noindent\textbf{Trang 21, §~4: Định lý hữu hạn cho lược đồ Picard.} --- Các
vấn đề hữu hạn thuộc loại được nêu trong mục này hầu như đã được giải quyết
hoàn toàn kể từ khi báo cáo hiện tại được biên soạn. Chúng tôi nêu các sự kiện
chính hiện đã biết theo hướng này. Để đơn giản hóa các phát biểu, chúng tôi
ngầm giả sử mọi tiền lược đồ Picard xuất hiện trong các phát biểu đều tồn tại,
mặc dù một sửa đổi hiển nhiên của các phát biểu cho phép loại bỏ mọi giả
thiết tồn tại. Trong phần sau, $S$ chỉ một lược đồ Noether, còn $X,Y$ là các
lược đồ proper trên $S$.

\medskip
\noindent(i) Cho $f:X\longrightarrow Y$ là một cấu xạ toàn ánh. Khi đó
\[
f^*: \operatorname{Pic}_{Y/S}\longrightarrow\operatorname{Pic}_{X/S}
\]
là một cấu xạ kiểu hữu hạn.

\src{304}
\noindent(ii) Giả sử $Y$ xạ ảnh trên $S$, được trang bị một môđun khả nghịch
ample tương đối trên $S$, và cho $X$ là lược đồ không của một tiết diện tùy ý
của môđun này, $f:X\longrightarrow Y$ là phép nhúng chính tắc. Cuối cùng, giả
sử các thành phần bất khả quy của các thớ của $Y/S$ có chiều $\geq3$. Khi đó
\[
f^*: \operatorname{Pic}_{Y/S}\longrightarrow\operatorname{Pic}_{X/S}
\]
là kiểu hữu hạn.

\noindent(iii) Giả sử $X$ xạ ảnh trên $S$ và mọi thớ hình học của nó đều
nguyên, có chiều $n$. Cho $\mathcal O_X(1)$ là một môđun khả nghịch ample trên
$X$, cho phép xác định các đa thức Hilbert. Cho $M$ là một tập con của
$\operatorname{Pic}_{X/S}$; khi đó $M$ tựa compact khi và chỉ khi, trong các
đa thức Hilbert
\[
a_0x^n+a_1x^{n-1}+a_2x^{n-2}+\cdots
\]
của các phần tử thuộc $M$, các hệ số $a_1$ và $a_2$ vẫn bị chặn, bỏ qua mọi hệ số có đơn thức mang bậc âm; do đó chỉ $a_1$ xuất hiện khi $n=1$, và cả hai hệ số đều không xuất hiện khi $n=0$.

\noindent(iv) Với mọi số nguyên $n\neq0$, đồng cấu lũy thừa bậc $n$ trong
tiền lược đồ nhóm $\operatorname{Pic}_{X/S}$ là một cấu xạ kiểu hữu hạn.

Lưu ý rằng (i) và (ii) cũng có nghĩa là (dưới các giả thiết tương ứng trong
hai phát biểu ấy) một tập con $M$ của $\operatorname{Pic}_{Y/S}$ tựa compact
khi và chỉ khi ảnh của nó trong $\operatorname{Pic}_{X/S}$ tựa compact. Suy
ra một môđun khả nghịch $\mathcal L$ trên $Y$ là $\tau$-tương đương với không
khi và chỉ khi ảnh ngược của nó trên $X$ có tính chất ấy; nói cách khác,
$\operatorname{Pic}_{Y/S}^\tau$ là ảnh ngược của
$\operatorname{Pic}_{X/S}^\tau$. Đặc biệt, để chứng minh tiền lược đồ thứ
nhất thuộc kiểu hữu hạn trên $S$, chỉ cần chứng minh điều đó cho tiền lược đồ
thứ hai, vì $\operatorname{Pic}_{Y/S}\longrightarrow
\operatorname{Pic}_{X/S}$ thuộc kiểu hữu hạn. Khi đó, dùng (i), Bổ đề Chow và
(iii), ta thu được:

\medskip
\noindent(v) $\operatorname{Pic}_{X/S}^\tau$ thuộc kiểu hữu hạn trên $S$.

Tổng quát hơn, việc kết hợp (i) cho một cấu xạ hữu hạn với (ii) cho phép quy
hầu hết các vấn đề hữu hạn về trường hợp $X/S$ có các thớ hình học nguyên và
chuẩn, chiều $\leq2$; thường thì, bằng cách áp dụng (i) cho một cấu xạ toàn
ánh không nhất thiết hữu hạn và dùng phép giải kỳ dị của các mặt đại số (do
Abhyankar thực hiện trong đặc số tùy ý), ta thậm chí quy được về trường hợp
$X/S$ trơn, do đó có các thớ hình học không kỳ dị chiều 2. Kết hợp điều này
với (v) và bất đẳng thức Picard--Igusa chặn hạng của nhóm Néron--Severi của
một mặt xạ ảnh không kỳ dị, ta chứng minh được sự khái quát hóa sau đây của
định lý hữu hạn Néron:

\medskip
\noindent(vi) Cho $X/S$ proper trên $S$ và ngoài ra tùy ý. Khi đó các nhóm
Néron--Severi
\[
\operatorname{Pic}_{X_i/k_i}/\operatorname{Pic}_{X_i/k_i}^\circ
\]
của các thớ hình học $X_i/k_i$ của $X/S$

\src{305}
đều thuộc kiểu hữu hạn, còn hạng của chúng và cấp của các nhóm con xoắn của
chúng vẫn bị chặn trên.

Cùng một phương pháp quy về trường hợp các mặt không kỳ dị, kết hợp với các
định lý đã biết trong trường hợp đó (cụ thể là định lý hữu hạn Néron và tính
không suy biến của dạng giao trên nhóm Néron--Severi), kéo theo:

\medskip
\noindent(vii) Cho $X$ là một lược đồ proper trên một trường đóng đại số. Khi
đó tồn tại hữu hạn đường cong đóng nguyên $C_i$ ($1\leq i\leq r$) trong $X$
có tính chất sau: với mọi tập con $M$ của $\operatorname{Pic}_{X/k}$, $M$ tựa
compact khi và chỉ khi các số nguyên
\[
\deg_{C_i'}\mathcal L\qquad(\mathcal L\in M)
\]
vẫn bị chặn (trong đó $C_i'$ chỉ chuẩn hóa của $C_i$).

Ở đây có thể lấy $r$ bằng hạng của nhóm Néron--Severi. Một khi đã biết nhóm
này thuộc kiểu hữu hạn, (vii) quy về sự kiện rằng các dạng tuyến tính trên
nhóm Néron--Severi được xác định bởi các đường cong $C$ trong $X$ chỉ đồng
thời triệt tiêu trên các phần tử xoắn của nhóm Néron--Severi. Trong trường hợp
$X$ xạ ảnh và không kỳ dị, kết quả này, cũng như (v), là do Matsusaka. Dùng
(vii), ta dễ dàng thu được đặc trưng sau đây của các môđun khả nghịch
$\tau$-tương đương với không trên $X$:

\medskip
\noindent(viii) Cho $X/k$ là một lược đồ proper trên một trường và
$\mathcal L$ là một môđun khả nghịch trên $X$. Các điều kiện sau là tương
đương:

\noindent a. $\mathcal L$ là $\tau$-tương đương với không;

\noindent b. Với mọi môđun coherent $F$, ta có
\[
\chi(F\otimes\mathcal L)=\chi(F),
\]
trong đó $\chi$ chỉ đặc trưng Euler--Poincaré;

\noindent b'. Như (b), nhưng với $F=\mathcal O_Y$, trong đó $Y$ là một lược
đồ con đóng nguyên chiều 1 trong $X$;

\noindent c. Với mọi $Y$ như trên, ký hiệu $Y'$ là đường cong chuẩn hóa, ta
có
\[
\deg\mathcal L_{Y'}=0.
\]

Nếu $X/k$ xạ ảnh và được trang bị một môđun khả nghịch ample
$\mathcal O_X(1)$, thì các điều kiện trước cũng tương đương với các điều kiện
sau:

\noindent d. Với mọi số nguyên $m$, $\mathcal L^{\otimes m}(1)$ là ample.

\noindent e. (Nếu $X$ nguyên.) Với mọi cặp số nguyên $m,n$, ta có
\[
\chi(\mathcal L^{\otimes m}(n))=\chi(\mathcal O_X(n))
\]

\src{306}
(nghĩa là (b) đúng khi lấy $F=\mathcal L^{\otimes m}(n)$).

Để chứng minh tính đủ của điều kiện cuối này, lưu ý rằng nó có nghĩa là các
đa thức Hilbert của $\mathcal L^{\otimes m}$ đều bằng nhau; do đó, theo tiêu
chuẩn Mumford (iii), các $\mathcal L^{\otimes m}$ vẫn nằm trong một tập tựa
compact của $\operatorname{Pic}_{X/k}$, nghĩa là ta có (a). Các điều kiện (b),
(b') và (c) phải được xem là những biến thể (trên một lược đồ proper tùy ý)
của khái niệm tương đương số, vốn thường được định nghĩa trên các đa tạp xạ
ảnh không kỳ dị. Đối với các đa tạp ấy, sự tương đương giữa (a) và (c) dĩ
nhiên đã được biết rõ (Matsusaka).

Tiêu chuẩn (e) cũng kéo theo kết quả sau:

\medskip
\noindent(ix) Cho $f:X\longrightarrow S$ là một cấu xạ xạ ảnh và phẳng có
các thớ hình học nguyên. Khi đó $\operatorname{Pic}_{X/S}^\tau$ vừa mở vừa
đóng trong $\operatorname{Pic}_{X/S}$.

Chúng tôi chỉ giới hạn ở một vài bình luận về chứng minh của các kết quả then
chốt (i), (ii), (iii) (kết quả (iv) hơi tách biệt với các kết quả khác và
được chứng minh chỉ bằng cách dùng (i) cho các cấu xạ hữu hạn, toàn ánh và
radiciel, cụ thể là cho một cấu xạ Frobenius). Đối với (i), ta sử dụng một
cách thiết yếu các ý tưởng của descent không phẳng (xem đặc biệt TDTE I,
trang 9). Do chỉ nhắm đến các kết quả hữu hạn, việc thiếu các tiêu chuẩn hữu
hiệu cho dữ liệu descent là vô hại. Mặt khác, Mumford gần đây đã chứng minh
một dạng hơi yếu hơn của (iii), cụ thể là tiêu chuẩn có sự tham gia của mọi
hệ số trong đa thức Hilbert. Lập luận của ông cực kỳ đơn giản và được gợi ý từ
chứng minh của một tiêu chuẩn ample của Nakai (do Nakai phát biểu cho các mặt
không kỳ dị và được Mumford khái quát hóa cho các cấu xạ xạ ảnh tùy ý). Tuy
nhiên, theo tôi, lập luận này chỉ đúng khi đặt thêm một hạn chế nhẹ lên các thớ
của $X/S$, tức điều kiện $(S_2)$ của Serre, vốn được thỏa mãn nếu các thớ hình
học là chuẩn. Sau đó ta dùng tiêu chuẩn bị hạn chế này trong chứng minh (ii):
tiêu chuẩn (i) thực sự cho phép quy về trường hợp $Y/S$ phẳng, có các thớ hình
học nguyên và chuẩn; rồi áp dụng tiêu chuẩn Mumford, ta dễ dàng quy về trường
hợp $X/S$ thỏa mãn các điều kiện tương tự. Từ giả thiết về chiều suy ra các
thớ hình học của $Y$ và $X$ có độ sâu $\geq2$ tại các điểm đóng của chúng;
điều này cho phép áp dụng các ``tiêu chuẩn tương đương'' dưới dạng được nêu
trong SGA 1962 (Báo cáo XII) và hoàn tất chứng minh của (ii). Một khi đã có
(i) và (ii), không khó để loại bỏ khỏi tiêu chuẩn Mumford mọi giả thiết về
tính chuẩn của các thớ và cũng chứng minh tiêu chuẩn ấy dưới dạng mạnh hơn
được nêu trong (iii).

\src{307}
Cuối cùng, lưu ý rằng chứng minh của (i) và (ii) cũng cho thấy rằng, trong
trường hợp $S$ là phổ của một trường $k$, cấu xạ
\[
\operatorname{Pic}_{Y/k}\longrightarrow\operatorname{Pic}_{X/k}
\]
là afin (chứ không chỉ thuộc kiểu hữu hạn).

\par
